\chapter{Every integer has an address}
\chaptertag{Primes and the fundamental theorem of arithmetic}

Positive integers factor. But why should the factorization into primes be unique? The arithmetic step that makes uniqueness possible is modest enough to miss.

\begin{definition}
An integer $p>1$ is \emph{prime} if its only positive divisors are $1$ and $p$. An integer $n>1$ that is not prime is \emph{composite}.
\end{definition}

\section{Euclid's lemma}

\begin{lemma}[Euclid]
If $p$ is prime and $p\divides ab$, then $p\divides a$ or $p\divides b$.
\end{lemma}

\begin{proof}
Suppose $p\ndivides a$. Since the only positive divisors of $p$ are $1,p$, we have $\gcd(p,a)=1$. B\'ezout gives $px+ay=1$. Multiply by $b$:
\[
pbx+aby=b.
\]
Both terms on the left are divisible by $p$, so $p\divides b$.
\end{proof}

\begin{warningbox}
The word ``prime'' is essential. We have $6\divides 2\cdot3$, but $6$ divides neither factor.
\end{warningbox}

Applying the lemma repeatedly shows: if a prime divides $a_1a_2\cdots a_k$, it divides at least one factor.

\section{Existence and uniqueness}

\begin{theorem}[Fundamental theorem of arithmetic]
Every integer $n>1$ can be written as a product of primes. Apart from rearranging the factors, this product is unique.
\end{theorem}

\begin{proof}[Existence]
Use strong induction on $n$. If $n$ is prime, we are done. If it is composite, write $n=ab$ with $1<a,b<n$. By induction, both $a$ and $b$ are products of primes, so $n$ is too.
\end{proof}

\begin{proof}[Uniqueness]
Suppose
\[
p_1p_2\cdots p_r=q_1q_2\cdots q_s
\]
with all $p_i,q_j$ prime. Since $p_1$ divides the right side, Euclid's lemma says $p_1=q_j$ for some $j$. Cancel this common factor. Repeat. No factor can remain on only one side, so the lists agree up to order.
\end{proof}

\section{An exponent-parity proof of irrationality}

Unique factorization can rule out rational numbers without trying to compute
their decimal expansions.

\begin{theorem}
The number $\sqrt2$ is irrational.
\end{theorem}

\begin{proof}
Suppose instead that $\sqrt2=a/b$ for integers $a,b$ with $b\ne0$. Squaring
and clearing the denominator gives
\[
  a^2=2b^2.
\]
By the fundamental theorem of arithmetic, the exponent of $2$ in the prime
factorization of $a^2$ is even: every exponent occurring in $a$ has been
doubled. The exponent of $2$ in $2b^2$ is odd: it is one plus twice the
exponent occurring in $b$. Uniqueness says that equal positive integers must
have the same exponent of every prime, so this is impossible.
\end{proof}

\begin{center}
\begin{tikzpicture}[
  node distance=12mm,
  parity/.style={draw=navy!55,rounded corners=2mm,fill=white,minimum width=48mm,
    minimum height=13mm,align=center},
  arr/.style={-{Stealth[length=2.2mm]},thick,coral}
]
  \node[parity] (left) {$a^2$ has an \textbf{even}\\number of factors $2$};
  \node[parity,right=of left] (right) {$2b^2$ has an \textbf{odd}\\number of factors $2$};
  \draw[arr] (left)--node[above,font=\small\bfseries] {cannot be equal} (right);
\end{tikzpicture}

\small The contradiction is not merely that one side is even; it is that the
two sides contain incompatible numbers of factors $2$.
\end{center}

The same proof works with any prime in place of $2$. More generally, a rational
number can be a perfect square only when every prime exponent in its reduced
numerator and denominator is even.

\begin{labbox}[title={Aside: existence without a witness}]
Here is a famous one-line use of the law of excluded middle. Set
\[
  x=(\sqrt2)^{\sqrt2}.
\]
Either $x$ is rational or it is irrational. In the first case, the two
irrational numbers $a=b=\sqrt2$ satisfy $a^b=x\in\mathbb Q$. In the second
case, the two irrational numbers $a=x$ and $b=\sqrt2$ satisfy
\[
  a^b=\left((\sqrt2)^{\sqrt2}\right)^{\sqrt2}=2.
\]
Thus, classically, there exist irrational $a$ and $b$ for which $a^b$ is
rational.

The proof is valid in classical mathematics, but it is not a constructive
proof: it uses the disjunction ``$x$ is rational or $x$ is not rational''
without deciding which alternative holds, so it does not identify which of
the two proposed pairs is a witness. Constructive and intuitionistic
mathematics do not in general accept this unrestricted use of excluded
middle.

Unique factorization gives an elementary explicit witness. Take
\[
  a=\sqrt2,
  \qquad
  b=\log_{\sqrt2}3.
\]
Then $a^b=3$. Moreover $b$ is irrational: if $b=p/q$ with $q>0$, then
raising $(\sqrt2)^{p/q}=3$ to the power $2q$ would give
\[
  2^p=3^{2q},
\]
contradicting unique prime factorization. This time the proof supplies the
pair that works. Much more deeply, the Gelfond--Schneider theorem implies that
$(\sqrt2)^{\sqrt2}$ is transcendental, so it retrospectively selects the
second branch of the classical argument.
\end{labbox}

Collecting repeated primes gives a coordinate system
\[
n=2^{e_2}3^{e_3}5^{e_5}\cdots,
\]
where only finitely many exponents are nonzero. Multiplication adds exponent vectors; gcd takes coordinatewise minima; lcm takes coordinatewise maxima.

\begin{center}
\begin{tikzpicture}[x=1.2cm,y=1.0cm]
  \draw[-{Stealth},navy] (0,0)--(4.5,0) node[right] {$e_2$};
  \draw[-{Stealth},navy] (0,0)--(0,3.7) node[above] {$e_3$};
  \foreach \x in {0,...,4} \foreach \y in {0,...,3} \fill[navy!25] (\x,\y) circle (1.3pt);
  \fill[teal] (3,2) circle (4pt) node[above right] {$72=2^3 3^2$};
  \fill[coral] (2,3) circle (4pt) node[above right] {$108=2^2 3^3$};
  \fill[gold!80!black] (2,2) circle (4pt) node[below left] {$\gcd=36$};
  \draw[dashed,gold!80!black] (2,2)--(3,2);
  \draw[dashed,gold!80!black] (2,2)--(2,3);
  \fill[navy] (3,3) circle (4pt) node[above right] {$\operatorname{lcm}=216$};
\end{tikzpicture}

\small Prime exponents turn gcd and lcm into geometry.
\end{center}

\begin{proposition}
If $n=\prod p^{e_p}$ and $m=\prod p^{f_p}$, then
\[
\gcd(n,m)=\prod p^{\min(e_p,f_p)},\qquad
\operatorname{lcm}(n,m)=\prod p^{\max(e_p,f_p)}.
\]
In particular, $\gcd(n,m)\operatorname{lcm}(n,m)=nm$.
\end{proposition}

\section{The sieve of Eratosthenes}

Unique factorization also explains a practical way to find every prime up to a
chosen bound $N$. Write down the integers $2,3,\dots,N$. Circle the first
uncrossed number, then cross out its larger multiples. Repeat. The circled
numbers are precisely the primes.

\begin{center}
\newcommand{\sieveprime}[1]{\tikz[baseline=(n.base)]\node[
  draw=teal,fill=teal!13,rounded corners=1mm,minimum size=6.5mm,
  inner sep=0pt,font=\small\bfseries,text=teal!70!black] (n) {#1};}
\newcommand{\sievegone}[1]{\tikz[baseline=(n.base)]{
  \node[minimum size=6.5mm,inner sep=0pt,font=\small,text=muted] (n) {#1};
  \draw[coral!65,line width=.45pt] (n.south west)--(n.north east);}}
\newcommand{\sieveone}[1]{\tikz[baseline=(n.base)]\node[
  minimum size=6.5mm,inner sep=0pt,font=\small,text=muted!70] (n) {#1};}
\renewcommand{\arraystretch}{1.25}
\begin{tabular}{*{10}{c}}
\sieveone{1}&\sieveprime{2}&\sieveprime{3}&\sievegone{4}&\sieveprime{5}&\sievegone{6}&\sieveprime{7}&\sievegone{8}&\sievegone{9}&\sievegone{10}\\
\sieveprime{11}&\sievegone{12}&\sieveprime{13}&\sievegone{14}&\sievegone{15}&\sievegone{16}&\sieveprime{17}&\sievegone{18}&\sieveprime{19}&\sievegone{20}\\
\sievegone{21}&\sievegone{22}&\sieveprime{23}&\sievegone{24}&\sievegone{25}&\sievegone{26}&\sievegone{27}&\sievegone{28}&\sieveprime{29}&\sievegone{30}\\
\sieveprime{31}&\sievegone{32}&\sievegone{33}&\sievegone{34}&\sievegone{35}&\sievegone{36}&\sieveprime{37}&\sievegone{38}&\sievegone{39}&\sievegone{40}\\
\sieveprime{41}&\sievegone{42}&\sieveprime{43}&\sievegone{44}&\sievegone{45}&\sievegone{46}&\sieveprime{47}&\sievegone{48}&\sievegone{49}&\sievegone{50}
\end{tabular}

\small Sieve to $50$: the boxed survivors are prime; every composite has been crossed out.
\end{center}

\begin{proposition}[Why the sieve stops early]
If $n>1$ is composite, then it has a prime divisor at most $\sqrt n$.
Consequently, to find all primes up to $N$, it is enough to cross out multiples
of primes $p\le\sqrt N$.
\end{proposition}

\begin{proof}
Write $n=ab$ with $1<a\le b$. Then $a\le\sqrt n$. A prime factor $p$ of $a$
therefore satisfies $p\le a\le\sqrt n$. Thus if an integer $n\le N$ survives
all sieving primes through $\sqrt N$, it cannot be composite.
\end{proof}

\begin{checkpoint}
For the display above, only the multiples of $2,3,5,$ and $7$ need to be
crossed out. The next prime is $11>\sqrt{50}$. Notice also that one may begin
crossing out at $p^2$: every smaller multiple of $p$ has already been handled
by an earlier prime.
\end{checkpoint}

\section{There is no last prime}

\begin{theorem}[Euclid]
There are infinitely many primes.
\end{theorem}

\begin{proof}
If $p_1,\dots,p_k$ were all primes, consider $N=p_1\cdots p_k+1$. Some prime divides $N$, but none of the listed primes does: division by any $p_i$ leaves remainder $1$. Contradiction.
\end{proof}

This proof does not say that $N$ is prime. It says that $N$ has a prime factor missing from the proposed complete list.

\begin{exercisebox}
\begin{enumerate}
  \item Factor $7560$ and use the exponent vectors to find the number of its positive divisors. \exlevel{Fluency}
  \item Compute $\gcd(2^7 3^2 11,2^4 3^5 5)$ and the lcm without multiplying first. \exlevel{Fluency}
  \item Prove that if $\gcd(a,b)=1$ and $a\divides c$, $b\divides c$, then $ab\divides c$. \exlevel{Reasoning}
  \item Reproduce the exponent-parity proof that $\sqrt2$ is irrational, then prove that $\sqrt p$ is irrational for every prime $p$. State exactly where unique factorization is used. \exlevel{Reasoning}
  \item Prove that every composite $n>1$ has a prime divisor at most $\sqrt n$. Turn this into a primality test. \exlevel{Reasoning}
  \item Use the sieve of Eratosthenes to list every prime at most $100$. Explain why, after sieving by $2,3,5,$ and $7$, no further crossing-out step is needed. \exlevel{Fluency}
  \item Is $n!+1$ always prime? Find the first counterexample, then explain why Euclid's proof is unaffected. \exlevel{Exploration}
  \item Prove that infinitely many primes are congruent to $3$ modulo $4$ by adapting Euclid's idea to $4p_1\cdots p_k-1$. You may use the fact that a product of numbers congruent to $1$ modulo $4$ is again $1$ modulo $4$. \exlevel{Challenge}
  \item Verify directly that $a=\sqrt2$ and $b=\log_{\sqrt2}3$ are irrational and satisfy $a^b=3$. What does this explicit witness add to the excluded-middle proof? \exlevel{Reasoning}
\end{enumerate}
\end{exercisebox}
